% methodology.tex -- Batch Aggregation Paper (RU-V4)

\section{Methodology}

\subsection{Protocol Design}

The batch aggregation protocol consists of three subsystems: a persistent queue
backed by a write-ahead log, a configurable batch aggregator, and a
checkpoint-based crash recovery mechanism.

\subsubsection{Write-Ahead Log (WAL)}

Every incoming transaction is first appended to an on-disk WAL before being added
to the in-memory queue. Each WAL entry contains:

\begin{center}
\texttt{\{seq, timestamp, tx, checksum\}}
\end{center}

where \texttt{seq} is a monotonically increasing sequence number,
\texttt{timestamp} is the arrival time (used for FIFO ordering), \texttt{tx} is
the full transaction payload, and \texttt{checksum} is a SHA-256 integrity digest.
The WAL uses a JSON-lines append-only format, chosen for debuggability at the
target scale ($<$10K tx/min). Durability is configurable: per-entry fsync for
maximum safety or group commit (fsync per $N$ entries) for throughput.

\subsubsection{Batch Formation Strategies}

Three strategies are implemented and evaluated:

\begin{itemize}
\item \textbf{SIZE}: Form a batch when the queue reaches a size threshold $N$.
  Guarantees fixed-size batches optimal for circuit proving. Unbounded latency
  under low load.
\item \textbf{TIME}: Form a batch every $T$ milliseconds regardless of queue
  size. Bounded latency but variable batch sizes that may not match circuit
  capacity.
\item \textbf{HYBRID}: Form a batch when either threshold triggers
  ($|queue| \geq N$ OR $elapsed \geq T$ with $|queue| > 0$). This combines the
  bounded latency of TIME with the bounded batch size of SIZE, matching the
  approach used by all major production ZK-rollups~\cite{chaliasos2024analyzing}.
\end{itemize}

\subsubsection{Checkpoint Protocol}

The WAL checkpoint marks the high-water sequence number below which transactions
are considered ``committed.'' After crash recovery, only entries above the
checkpoint are replayed into the queue.

In our initial design, the checkpoint advanced at \emph{batch formation} time:
once transactions were dequeued into a batch, the WAL considered them committed.
Section~\ref{sec:formal} demonstrates why this is incorrect. The corrected
protocol defers the checkpoint to \emph{batch processing} time (after successful
ZK proof generation and L1 confirmation), ensuring that unprocessed batch
transactions remain recoverable from the WAL.

\subsubsection{Determinism Contract}

Batch identity is defined as $\text{BatchID} = \text{SHA-256}(h_1 \| h_2 \| \ldots \| h_n)$
where $h_i$ are the transaction hashes in FIFO order. Transactions within a
batch are ordered by arrival timestamp with sequence number tie-breaking.
This guarantees that the same set of transactions always produces an identical
batch, enabling replay-based verification.

\subsection{Empirical Evaluation}

\subsubsection{Benchmark Suite}

Five experimental phases evaluate the protocol across its parameter space:

\begin{enumerate}
\item \textbf{Strategy comparison}: SIZE vs.\ TIME vs.\ HYBRID under identical
  load conditions.
\item \textbf{Batch size sweep}: Threshold values $\{4, 8, 16, 32, 64\}$ to
  match ZK circuit batch capacities established in prior work (RU-V2).
\item \textbf{Arrival rate sweep}: Transaction rates $\{50, 100, 200, 500, 1000\}$
  tx/min spanning the target range.
\item \textbf{Time threshold sweep}: Values $\{1000, 2000, 5000\}$ ms to evaluate
  latency bounds under low load.
\item \textbf{fsync strategy comparison}: Per-entry fsync vs.\ group commit
  (batch size 16).
\end{enumerate}

Each configuration runs 30 replications with different random seeds. Two
warmup batches are discarded. Synthetic transactions are generated with
Poisson-like arrival jitter (exponential inter-arrival times around the mean
rate). We report mean, standard deviation, and 95\% confidence intervals.
CI half-width must be within 10\% of the mean (all configurations pass this
criterion).

\subsubsection{Determinism Tests}

450 test cases across 3 strategies and 5 batch sizes (30 replications each).
For each test, the same set of deterministic transactions is fed to the
aggregator twice, and the resulting batch IDs and transaction orderings are
compared for exact equality.

\subsubsection{Crash Recovery Tests}

150 test cases across 5 crash scenarios (30 replications each):
\begin{enumerate}
\item Crash after enqueue, before batch formation
\item Crash during batch formation (mid-batch)
\item Corrupted WAL entry (simulated partial write)
\item Crash after checkpoint (clean state)
\item Multiple sequential crashes (3 cycles)
\end{enumerate}

Each test verifies zero transaction loss after recovery.

\subsection{Formal Verification}
\label{sec:formal}

We formalize the batch aggregation protocol in TLA+ (Temporal Logic of
Actions)~\cite{lamport2002specifying} with the following structure:

\begin{itemize}
\item \textbf{8 state variables}: \texttt{queue}, \texttt{wal},
  \texttt{checkpointSeq}, \texttt{batches}, \texttt{processed}, \texttt{pending},
  \texttt{systemUp}, \texttt{timerExpired}
\item \textbf{6 actions}: \texttt{Enqueue(tx)}, \texttt{FormBatch},
  \texttt{ProcessBatch}, \texttt{Crash}, \texttt{Recover}, \texttt{TimerTick}
\item \textbf{6 safety invariants}: \texttt{TypeOK}, \texttt{NoLoss},
  \texttt{NoDuplication}, \texttt{QueueWalConsistency}, \texttt{FIFOOrdering},
  \texttt{BatchSizeBound}
\item \textbf{1 liveness property}: \texttt{EventualProcessing} (every tx
  eventually appears in a processed batch)
\end{itemize}

The crash model clears all volatile state (queue, batches) while preserving
durable state (WAL, checkpoint, processed). Recovery replays the WAL from the
checkpoint to reconstruct the queue. The nondeterministic \texttt{TimerTick}
over-approximates real-time behavior (fires at any point), making safety
results strictly stronger than reality.

Model checking uses TLC with exhaustive breadth-first search. The initial
model (v0) uses 10 transaction identifiers and batch threshold 4. The
corrected model (v1-fix) uses 4 identifiers and threshold 2 for complete
state-space exploration.
